\section{Аппроксимация функциональных зависимостей}
\label{sec:function-approximation}

\subsection{Постановка задачи}

Аппроксимацией называется замена сложной функции $f$ более простой функцией
$\varphi$ из заранее выбранного класса $V$:
\[
    f(x)\approx\varphi(x),
    \qquad \varphi\in V.
\]
Исходная зависимость может быть задана аналитически, таблицей значений
$(x_i,y_i)$ или результатами эксперимента. Класс $V$ выбирают с учётом требуемой
точности, гладкости, вычислительной стоимости и физического смысла модели.

Следует различать:
\begin{itemize}
    \item интерполяцию: $\varphi(x_i)=y_i$ во всех узлах;
    \item аппроксимацию в среднем: минимизируется интегральная или дискретная
          ошибка;
    \item равномерное приближение: минимизируется максимальная ошибка;
    \item регрессию: данные считаются зашумлёнными, точное прохождение через них
          не требуется.
\end{itemize}

\subsection{Нормы и критерии качества}

Качество приближения оценивается нормой ошибки $e=f-\varphi$. На отрезке $[a,b]$
часто используют
\[
\|e\|_{\infty}=\max_{x\in[a,b]}|e(x)|,
\qquad
\|e\|_{2}=\left(\int_a^b |e(x)|^2\,dx\right)^{1/2}.
\]
Для табличных данных:
\[
    \|e\|_{2,h}=
    \left(\sum_{i=1}^{N}w_i|y_i-\varphi(x_i)|^2\right)^{1/2},
\]
где $w_i>0$ --- веса наблюдений.

Абсолютная ошибка измеряется в единицах величины, относительная --- отношением к
масштабу функции. Для данных с различной дисперсией разумно использовать веса,
обратно пропорциональные дисперсии ошибок.

\subsection{Линейная аппроксимационная модель}

Пусть
\[
    \varphi(x;c)=\sum_{j=0}^{m}c_j\phi_j(x),
\]
где базисные функции $\phi_j$ заданы, а коэффициенты $c_j$ неизвестны. Такая
модель линейна по параметрам, даже если функции $\phi_j$ нелинейны по $x$.

В дискретном случае вводят матрицу дизайна
\[
    A_{ij}=\phi_j(x_i),
    \qquad
    Ac\approx y.
\]
При квадратичном критерии задача сводится к методу наименьших квадратов:
\[
    \min_c\|Ac-y\|_2^2.
\]
Геометрически решение является ортогональной проекцией вектора данных на
столбцовое пространство $A$.

\subsection{Основные классы аппроксимирующих функций}

\begin{enumerate}
    \item \textbf{Алгебраические многочлены}
    \[
        \varphi(x)=c_0+c_1x+\cdots+c_mx^m.
    \]
    Они просты, но степенной базис может быть плохо обусловлен при больших $m$.

    \item \textbf{Ортогональные разложения}
    \[
        \varphi_m(x)=\sum_{j=0}^{m}c_j\phi_j(x),
    \]
    где $\phi_j$ ортогональны в выбранном скалярном произведении. Примеры:
    многочлены Лежандра и Чебышева, тригонометрическая система.

    \item \textbf{Сплайны} --- кусочно-полиномиальные функции с заданной
    гладкостью. Они устойчивее глобальных многочленов и локально реагируют на
    изменение данных.

    \item \textbf{Рациональные функции}
    \[
        \varphi(x)=\frac{P_m(x)}{Q_n(x)}.
    \]
    Они хорошо описывают функции с полюсами и насыщением, но задача подбора
    коэффициентов обычно нелинейна.

    \item \textbf{Экспоненциальные и степенные модели}:
    \[
        y=ae^{bx},\qquad y=ax^b.
    \]
    Иногда их линеаризуют логарифмированием, однако это изменяет структуру ошибки.
\end{enumerate}

\subsection{Теорема Вейерштрасса}

Если $f$ непрерывна на $[a,b]$, то для любого $\varepsilon>0$ существует
многочлен $P$ такой, что
\[
    \max_{x\in[a,b]}|f(x)-P(x)|<\varepsilon.
\]
Теорема гарантирует принципиальную возможность равномерного полиномиального
приближения, но не утверждает, что многочлен высокой степени численно удобен.
На практике глобальная аппроксимация высокой степени может быть неустойчивой;
часто лучше применять сплайны или ортогональные базисы.

\subsection{Равномерное наилучшее приближение}

Задача Чебышёва имеет вид
\[
    \min_{\varphi\in V}\|f-\varphi\|_{\infty}.
\]
Для многочленов степени не выше $m$ наилучшее равномерное приближение
характеризуется чередованием ошибки: при стандартных условиях ошибка достигает
максимального по модулю значения с чередующимися знаками не менее чем в $m+2$
точках. На этом основан алгоритм Ремеза.

\subsection{Выбор сложности модели}

Увеличение числа параметров уменьшает ошибку на обучающих данных, но может
привести к переобучению и сильной чувствительности к шуму. Поэтому учитывают:
\begin{itemize}
    \item ошибку на независимой проверочной выборке;
    \item обусловленность матрицы задачи;
    \item физическую интерпретируемость коэффициентов;
    \item устойчивость к изменению данных;
    \item регуляризацию, например
    \[
        \min_c\bigl(\|Ac-y\|_2^2+\lambda\|Lc\|_2^2\bigr).
    \]
\end{itemize}

\subsection{Что важно сказать на экзамене}

Аппроксимация отличается от интерполяции отсутствием требования точного
совпадения во всех узлах. Нужно назвать критерий ошибки, выбрать класс функций и
объяснить получение коэффициентов. Хорошая модель должна быть не только точной,
но и устойчивой, умеренно сложной и адекватной природе данных.
